% Introduction

\section{Introduction}
\label{sec:introduction}

Runtime threat detection is a critical component of modern security infrastructure, tasked with identifying malicious activities as they occur within computing systems.
The proliferation of sophisticated attacks---from container escapes to advanced persistent threats (APTs)---demands detection systems that combine comprehensive visibility, intelligent analysis, and minimal performance impact.
However, existing approaches face fundamental trade-offs that limit their effectiveness.

\textbf{The Detection Trilemma.}
Current runtime security systems must choose among three competing properties:
(1) \textit{Detection Capability}: the ability to identify complex, previously unseen attacks through pattern analysis or machine learning;
(2) \textit{Performance}: minimal latency and resource overhead to avoid impacting production workloads;
and (3) \textit{Security Isolation}: the ability to safely execute third-party detection rules without risking host compromise.
As shown in Table~\ref{tab:trilemma}, no existing system achieves all three properties simultaneously.

\begin{table}[t]
\centering
\caption{The Detection Trilemma: Existing systems sacrifice at least one property.}
\label{tab:trilemma}
\begin{tabular}{lccc}
\toprule
\textbf{System} & \textbf{Detection} & \textbf{Performance} & \textbf{Isolation} \\
\midrule
Falco (rules) & Low & Medium & Low \\
ML in userspace & High & Low & Medium \\
Kernel modules & Medium & High & None \\
\textsc{WasmSentinel} & \textbf{High} & \textbf{High} & \textbf{High} \\
\bottomrule
\end{tabular}
\end{table}

Rule-based systems like Falco~\cite{falco} and Tetragon~\cite{tetragon} achieve reasonable performance by matching events against predefined patterns, but their detection capability is limited to known attack signatures.
Machine learning approaches offer superior detection of novel attacks but typically incur prohibitive overhead---Python-based ML pipelines can add hundreds of milliseconds of latency per event.
Kernel-level solutions provide excellent performance but lack isolation, creating risks when executing community-contributed detection logic.

\textbf{Our Insight.}
We observe that the perceived incompatibility between ML detection and performance stems from implementation choices rather than fundamental limitations.
By carefully designing the detection pipeline---from data collection through inference---we can achieve ML-grade detection capability with microsecond-level latency.
The key enablers are:
(1) \textit{WebAssembly's near-native execution speed} with strong sandboxing guarantees;
(2) \textit{lightweight neural network architectures} specifically designed for the threat detection domain;
and (3) \textit{zero-copy data transfer} between the kernel and Wasm runtime.

\textbf{WasmSentinel Overview.}
We present \textsc{WasmSentinel}, a high-performance runtime threat detection framework that executes machine learning models within WebAssembly sandboxes.
\textsc{WasmSentinel} introduces a novel \textit{hybrid detection architecture} that combines the efficiency of STIDE-based sequence analysis~\cite{stide} with the accuracy of neural network inference.
The fast path uses n-gram hashing for immediate anomaly detection, while suspicious events are escalated to a lightweight MLP for deep analysis---all within a single Wasm module.

\textbf{Contributions.}
This paper makes the following contributions:

\begin{itemize}
    \item \textbf{WasmML:} We present the first system to run ML inference inside a Wasm sandbox for runtime threat detection, achieving 16.5$\mu$s mean inference latency with a 180KB portable binary. (Section~\ref{sec:design-wasmml})

    \item \textbf{Hybrid Detection Architecture:} We design a two-tier detection strategy combining STIDE n-gram analysis with lightweight MLP, achieving F1=0.72 on the ADFA-LD benchmark while maintaining sub-millisecond latency. (Section~\ref{sec:design-hybrid})

    \item \textbf{Zero-Copy Data Channel:} We implement an efficient eBPF-to-Wasm data path that eliminates JSON serialization, reducing per-event overhead by 10$\times$ compared to conventional approaches. (Section~\ref{sec:design-zerocopy})

    \item \textbf{Comprehensive Evaluation:} We demonstrate that \textsc{WasmSentinel} outperforms Falco by 485$\times$, Tetragon by 121$\times$, and Sysdig by 303$\times$ in inference latency while achieving comparable or better detection accuracy. (Section~\ref{sec:evaluation})
\end{itemize}

\textbf{Paper Organization.}
Section~\ref{sec:background} provides background on eBPF, WebAssembly, and ML-based detection.
Section~\ref{sec:threat-model} defines our threat model.
Section~\ref{sec:design} presents the system design, and Section~\ref{sec:implementation} describes implementation details.
Section~\ref{sec:evaluation} reports our experimental evaluation.
Section~\ref{sec:discussion} discusses limitations and future work.
Section~\ref{sec:related-work} surveys related work, and Section~\ref{sec:conclusion} concludes.
